\documentclass{article}

% Language setting
% Replace `english' with e.g. `spanish' to change the document language
\usepackage[english]{babel}

% Set page size and margins
% Replace `letterpaper' with`a4paper' for UK/EU standard size
\usepackage[letterpaper,top=0.5cm,bottom=0.5cm,left=1cm,right=1cm,marginparwidth=1.75cm]{geometry}

% Useful packages
\usepackage{amsmath}
\usepackage{graphicx}
\usepackage{xcolor}
\usepackage[colorlinks=true, allcolors=blue]{hyperref}

\usepackage{titlesec}
\titleformat*{\section}{\bfseries}
\font\titlefont=cmr12 at 12pt
\font\authorfont=cmr12 at 10pt
\title{{\titlefont Comments, Errata, and Addenda For My Articles}}
\author{{\authorfont Ethan Goan}}
\date{\small \today}
% Redefinition of ToC command to get centered heading
\usepackage{tocloft, sectsty}
\sectionfont{\fontsize{10}{12}\selectfont}

\renewcommand{\sectionfont}{\smallfont\bfseries}
\begin{document}
\maketitle
\tableofcontents

% NOTE TO MYSELF
% Don't move or change the name of this, just leave here and then recompile as needed

\section{Uncertainty in Real-Time Semantic Segmentation on Embedded Systems}
Corrected on arXiv

\textbf{Eq.6}  $ p(\mathbf{w} | \mathbf{X}, \mathbf{Y} , \theta ) \propto \Big[\prod_{i=0}^{N-1}\mathcal{N}\Big({\color{red}\mathbf{y}}|f(\mathbf{x}_{i};\theta) \mathbf{w}, \sigma^{2}\Big)\Big] p(\mathbf{w}).$

\textbf{Eq.10 and 12} should read $\sigma^{2}_{N}(\hat{\mathbf{x}}) =
\sigma^{2} + {\color{red} \Phi}(\hat{\mathbf{x}})^{T}\Sigma_{\pi}{\color{red} \Phi}(\hat{\mathbf{x}})$ and
$\sigma^{2}_{N}(\hat{\mathbf{x}}) \approx \sigma^{2} +
{\color{red}\Phi}(\hat{\mathbf{x}})^{T}\Sigma_{\text{diag}}{\color{red}\Phi}(\hat{\mathbf{x}})$

\textbf{Eq. 15}: The approximation from ref [34] in the paper provides the Taylor Series approximation of

$ \mathcal{N}\Big(\frac{\mu_{j}}{\mu_{d}},\frac{\mu_{j}^{2}}{\mu_{d}^{2}}\big(\sigma^{2}_{j} / \mu_{j}^{2} +  \sigma^{2}_{d} / \mu_{d}^{2} \big)\Big)$. The issue with this is the square of the $\mu_{d}$ terms, as this is at risk of numerical overflow in single-point precision. Instead we use $\sigma^{2} \approx \mu_{j} / \mu_{d}\sqrt{\sigma^{2}_{j} + \sigma^{2}_{d}}$, which emperically follwed a similar trend and was less of a risk for numerical overflow (though not immune).

In Table 1, predictive performance metrics for PIDNet used the pre-trained weights from the official repo for the testing set of CityScapes, which appears to have included the validation data during training. The mIOU for PIDNet and Bayes-PIDNet with the weights for the validation set are 0.761 and 0.758 respectively.

\section{Piecewise Deterministic Markov Processes for Bayesian Neural Networks}
Corrected on arXiv

\textbf{Eq. 11}: $t_{i} =  -b_{i} / a_{i} + \sqrt{b_{i}^{2} +  a_{i}^{2}t_{i-1}^{2} + 2a_{i}b_{i}t_{i-1} - 2a_{i}log(1 - U)}\big ) /a_{i}$

\textbf{Bug}: When preparing my thesis, I found a bug in the code where gradients were not being appropriately scaled by the mini-batch and dataset size. Experiments rerun with results on arXiv. I also added some clarification to the event rate sampling algorithm which clarifies some of the results of how thinning work. This is captured by the $R$ term in Algorithm 2, and we explore this more in Appendix A. Code updated on Github as well.

\section{Bayesian Neural Networks: An Introduction and Survey}
Corrected on arXiv

Typos in, \textbf{Eqns 29-31} for leap frog integration. Have changed to match version given by Neal which is simpler.

\textbf{Eq. 28} missing factor of $1/2$.

For the $\alpha$ divergence on page 31, the limits to relate to forward/reverse KL divergence and Hellinger distance were not correct. It should be, $KL(q|p): \alpha \rightarrow 0, KL(p|q): \alpha \rightarrow 1,$ Hellinger: $ \alpha \rightarrow 0.5$.
\end{document}
